\chapter{Conclusion and Future Work}

\section{Answers to the research
questions}

\textbf{RQ1 --- retrieval effectiveness against text-only baselines.} A
multi-modal pipeline retrieving on visual embeddings outperforms
transcript-based retrieval by an order of magnitude at moment
granularity on everyday video (corpus-scope R@10 .111 vs .011 at tIoU ≥
0.5). The fully developed system, a SigLIP index with LLM query
decomposition, nearly doubles that initial visual baseline in turn (R@1
.026 \ensuremath{\rightarrow} .048; video-scope tIoU@1 .230). Scope
analysis attributes the remaining difficulty to cross-video confusion
rather than within-video localisation: given the right video, the system
places a correct moment in its top five for two-thirds of questions at
the looser τ = 0.3 (R@5 .677). The answer to RQ1 is therefore
affirmative with a precise caveat: \emph{visual} retrieval is effective
and improvable, while retrieval on what is \emph{said} is, on this
domain, not a viable baseline so much as a cautionary one. §5.5 shows
that the caveat cuts both ways, since on speech-dense video the same
pipeline's transcript channel doubles its visual channel.

\textbf{RQ2 --- the multi-modal embedding strategy.} The experiments
support a strategy, not merely a ranking: index the visual channel with
the strongest affordable contrastive backbone (SigLIP SO400M dominated
both CLIP variants, +69\% relative R@1 as the index); if the strongest
backbone is unaffordable at corpus scale, index cheaply and re-rank its
top candidates with the stronger model, which matched the expensive
index to within measurement noise; translate questions into the
encoder's native register (scene captions) via LLM decomposition for
deep-rank recall; and do \emph{not} fuse in a weak text channel at any
weight, nor re-rank with a text cross-encoder, since both degrade
quality for an identified cause (sparse, multilingual, conversational
speech). For the audio channel, which enters this system as transcribed
speech (§1.2), the optimal strategy on the primary benchmark is the null
strategy; on-screen-text OCR was out of scope and remains untested. The
speech-dense replication (§5.5) then completes the answer. The ordering
is domain-conditional by measurement rather than conjecture: on YouCook2
the transcript channel doubles the visual channel and late fusion
becomes the best configuration. The optimal embedding strategy is
therefore not a fixed recipe but a measured, domain-adaptive choice, and
this dissertation's harness is the instrument that measures it.

\textbf{RQ3 --- agentic decomposition and synthesis.} Query
decomposition, temporal tool use, and bounded self-reflection each
demonstrably contribute. At the retrieval layer, decomposition is
safe-by-construction (the original query is always retained) and lifts
recall. At the answering layer, a LangGraph state machine that anchors
events, walks the timeline, and audits its own citations raises five-way
QA accuracy from .447 to .547 over a single-loop agent, with the largest
gains on causal and temporal-current questions. The prior-only baselines
calibrate what that machinery buys: with text-only evidence even the
better agent only recovers the bare model's answer prior (.560), because
the channel's genuine gains (temporal-current) are cancelled by its
losses on silent visual actions (temporal-next, where delivered
transcripts drag the model below its own prior). The controlled
evidence-channel experiment completes the answer: agentic machinery
alone cannot compensate for evidence the model cannot perceive ---
supplying keyframes to the answerer nearly doubled temporal-next
accuracy (.341 \ensuremath{\rightarrow} .636) with everything else held
fixed, and the full multimodal stack is the only configuration to beat
its model's prior (+.167). An effective video-QA agent is therefore
\emph{jointly} a retrieval planner and a multimodal reader; either half
alone leaves the other's headroom unrealised.

\section{Future work}

Five directions follow directly from the limitations identified in
Chapter 6, ordered roughly by evidence-per-effort.

\textbf{Repair the text channel, then re-test fusion.} Machine-translate
transcripts to English (or adopt a multilingual text-retrieval encoder)
before indexing, disentangling the informativeness/readability confound
of §6.2 --- and only then revisit late fusion and text re-ranking, whose
negative results are conditional on the broken channel.

\textbf{Scale the best-configuration evaluation.} The
graph-plus-multimodal combination now exists and reaches .727 on the
150-question sample (§5.4.4); the direct next step is statistical rather
than architectural --- evaluating the full grounded split (approximately
US\$250) would shrink the confidence interval from ±7 to ±1.5 points and
settle the comparison with published zero-shot systems, and a multimodal
\emph{reflection} pass (auditing the keyframes themselves) is the one
untested component interaction.

\textbf{Complete the open-source model comparison.} The provider
abstraction already targets any OpenAI-compatible endpoint; running the
QA suite against local Qwen/Llama models (the plumbing and evaluation
scripts exist) would quantify the API-versus-open trade-off the project
plan anticipated.

\textbf{Adaptive segmentation, on a tested segmenter.} Fixed 8-second
windows impose a structural ceiling on tIoU (§6.2); shot-aligned or
query-adaptive windowing would raise the ceiling itself, and the
evaluation harness can measure the gain directly. Reworking the
segmenter is also the point at which the missing unit tests (§4.7)
should be written: the boundary arithmetic and the rank-fusion and
imputation logic are exactly the code whose errors surface as
plausible-looking metrics rather than crashes.

\textbf{Deepen the second domain, and run the user study.} The
speech-dense retrieval replication is done (§5.5); its natural
extensions are to repeat the QA-stage and agent experiments on YouCook2,
replicate the \emph{positive} results (decomposition, re-ranking,
backbone ordering) cross-domain, and derive a quantitative
domain-adaptive fusion policy from measured speech density. The descoped
user study remains the right instrument for claims about end-user
utility that no offline metric can support.

\section{Closing remark}

The project set out to make a moment of video as findable as a sentence
of text, under the discipline that every design choice be measured. Its
most useful legacy may be less the headline numbers than the shape of
the evidence: failures separated into recall, precision, and evidence
problems; negative results retained with their causes; and a pipeline in
which each remedy could be switched on, switched off, and priced. Video
understanding systems will continue to change quickly; an evaluation
harness that can say \emph{which part helped, by how much, and why} will
remain the transferable asset.
